\documentclass{article}
\usepackage{graphicx} % Required for inserting images
\usepackage{tikz}
\usepackage{physics}
\usepackage{gensymb}
\usepackage{amsmath, amsfonts, mathtools, amsthm, amssymb, mathrsfs}
\usepackage{titlesec}
\usepackage{fancyhdr}
\usepackage{hyperref}
\usepackage{float}
\usepackage{booktabs}
\usepackage{enumitem}
\usepackage{subcaption}
\usepackage{multicol}
\usepackage{import}
\usepackage{bm}
\usepackage[usenames,dvipsnames]{xcolor}
\usepackage[a4paper, margin=1in]{geometry}
\usepackage{calligra}

\DeclareMathAlphabet{\mathcalligra}{T1}{calligra}{m}{n}
\DeclareFontShape{T1}{calligra}{m}{n}{<->s*[2.2]callig15}{}
\newcommand{\scriptr}{\mathcalligra{r}\,}
\newcommand{\boldscriptr}{\pmb{\mathcalligra{r}}\,}
% From https://florianbuerzle.wordpress.com/2011/07/18/griffiths-electrodynamics-script-r/
% The configuration for boxes and parts of the overall idea come from Pingbang Hu: https://github.com/sleepymalc
% Use a relative path (forward slashes or "..") to avoid backslash-escapes
\subimport{../}{header.tex}

\title{\textbf{PHYS 217 E\&M 1}}
\author{Autumn Weaver}
\date{Fall 2026}

\begin{document}

\maketitle

\tableofcontents{\textbf{}}

\pagebreak

\pagestyle{fancy}
\fancyhead{}
\fancyhead[L]{\thetitle}
\fancyhead[R]{Autumn Weaver}

\section{Vector Analysis}
\subsection{Vector Algebra}
\begin{definition}
    A \textbf{vector} is a quantity that has both direction and magnitude.
\end{definition}

\begin{definition}
    A \textbf{scalar} is a quantity that only has magnitude.
\end{definition}

\subsubsection{Vector Operations}

There are four vector operations:

\textbf{Addition}:
\begin{itemize}
    \item \(\mathbf{A} +\mathbf{B} \) is the vector from the start of \(\mathbf{A} \) to the end of \(\mathbf{B} \)   
    \item Commutative: \(\mathbf{A} +\mathbf{B} =\mathbf{B} +\mathbf{A} \) 
    \item Associative: \((\mathbf{A} +\mathbf{B} )+\mathbf{C} =\mathbf{A} +(\mathbf{B} +\mathbf{C} )\)
    \item Subtracting a vector is just adding its opposite: \(\mathbf{A} -\mathbf{B} =\mathbf{A} +(-\mathbf{B}) \) 
\end{itemize}


\textbf{Multiplication by Scalar}:
\begin{itemize}
    \item \(a \mathbf{A} \) is a vector with the same direction as \(\mathbf{A} \) but \(a\) times the magnitude
    \item Distributive: \(a(\mathbf{A} +\mathbf{B} )=a \mathbf{A} +a \mathbf{B} \)    
\end{itemize}


\textbf{Dot Product} (Inner Product or Scalar Product):
\begin{itemize}
    \item \(\mathbf{A} \cdot \mathbf{B} \equiv AB\cos{\theta}\) 
    \item Results in scalar
    \item Commutative: \(\mathbf{A} \cdot \mathbf{B} =\mathbf{B} \cdot \mathbf{A} \)
    \item Distributive: \(\mathbf{A} \cdot (\mathbf{B} +\mathbf{C} )=\mathbf{A} \cdot \mathbf{B} +\mathbf{A} \cdot \mathbf{C} \)
    \item Geometrically, \(\mathbf{A} \cdot \mathbf{B} \) is the product of \(A\) and the projection of \(\mathbf{A} \) along \(\mathbf{B} \)    
\end{itemize}

\textbf{Cross Product of Two Vectors}:
\begin{itemize}
    \item \(\mathbf{A} \cross  \mathbf{B} \equiv AB\sin{\theta}\hat{\mathbf{n}}\) where \(\mathbf{\hat{n}}\) is a unit vector perpendicular to the plane of \(\mathbf{A} \) and \(\mathbf{B} \). The direction 
    \item is determined by the \textbf{right hand rule}, where you curl your right hand from \(\mathbf{A} \) to \(\mathbf{B} \) and take the direction of your thumb
    \item Distributive: \(\mathbf{A} \cross (\mathbf{B}+\mathbf{C})=(\mathbf{A} \cross \mathbf{B} )+(\mathbf{A} \cross \mathbf{C} ) \)
    \item Not commutative: \(\mathbf{B} \cross \mathbf{A} =-(\mathbf{A} \cross \mathbf{B} )\)
    \item Result is a \textbf{pseudovector} and behaves like a vector under rigid transformations like rotations, but not under certain nonrigid transformations. Notably, \(-\mathbf{A} \cross -\mathbf{B} =\mathbf{A} \cross \mathbf{B} \)
    \item Geometrically, \(\lvert \mathbf{A} \cross \mathbf{B} \rvert \) is the area of the parallelogram with sides \(\mathbf{A} \) and \(\mathbf{B} \)  
\end{itemize}

When working in Cartesian coordinates, vectors can be expressed in terms of \textbf{basis vectors}:
\[
    \mathbf{A} =A_x \hat{\mathbf{x}}+A_y \mathbf{\hat{y}} +A_z \mathbf{\hat{z}}  
\]

with \(A_x\), \(A_y\), and \(A_z\) being the \textbf{components} of \(\mathbf{A} \).

Vector operations can be done with the components, but the reader should already know how to do this.

There are two types of triple vector products:

\textbf{Scalar Triple Product}:
\begin{itemize}
    \item \(\mathbf{A} \cdot (\mathbf{B}  \cross \mathbf{C} )=\mathbf{B} \cdot (\mathbf{C} \cross \mathbf{A} )=\mathbf{C} \cdot (\mathbf{A}  \cross \mathbf{B} )\)
    \item Volume of parallelepiped generated by \(\mathbf{A} \), \(\mathbf{B} \), and \(\mathbf{C} \)    
\end{itemize}

\textbf{Vector Triple Product}:
\begin{itemize}
    \item \textbf{BAC-CAB rule}: \(\mathbf{A} \cross (\mathbf{B} \cross \mathbf{C} )=\mathbf{B} (\mathbf{A} \cdot \mathbf{C} )-\mathbf{C} (\mathbf{A} \cdot \mathbf{B} )\) 
\end{itemize}

\subsubsection{Nomenclature}
\begin{definition}
    The \textbf{position vector} of a point is the vector to that point from the origin:
    \[
        \mathbf{r} \equiv x \mathbf{\hat{x}} +y \mathbf{\hat{y}} +z \mathbf{\hat{z}} 
    \]
\end{definition}

\begin{definition}
    The \textbf{infinitesimal displacement vector} from \((x,y,z)\) to \((x+dx, y+dy, z+dz)\) is:
    \[
        d \mathbf{I} =dx \mathbf{\hat{x}} + dy \mathbf{\hat{y}} +dz \mathbf{\hat{z}} 
    \]  
\end{definition}

\begin{definition}
    A \textbf{separation vector} is the vector from a \textbf{source point} to a \textbf{field point}:
    \[
        \boldscriptr \equiv \mathbf{r} - \mathbf{r^{\prime}}
    \]
\end{definition}

\subsubsection{Vector Transformations}

Coordination transformation from rotation about an arbitrary axis can be represented by the use of 
a \textbf{rotation matrix}:
\[
    \begin{pmatrix} \bar{A_x} \\ \bar{A_y} \\ \bar{A_z} \end{pmatrix} = \begin{pmatrix} R_{xx} & R_{xy} & R_{xz} \\ R_{yx} & R_{yy} & R_{yz} \\ R_{zx} & R_{zy} & R_{zz} \end{pmatrix} \begin{pmatrix} A_x \\ A_y \\ A_z \end{pmatrix}   
\]

or more concisely:
\[
    \bar{A_i}=\sum_{j=1}^3 R_{ij}A_j 
\]

\begin{definition}
    A transformation is \textbf{unitary} if it does not change the magnitude of the object on which it operates.
\end{definition}

If rotation matrix \(\overleftrightarrow{R}\) is unitary, then:
\[
    \sum_i^3R_ijR_ik=\delta_{jk}
\]

where:
\[
    \delta_{jk}=\begin{cases} 1 & \text{if } i=j \\ 0 & \text{if } i\neq j \end{cases}
\]

is the Kronecker delta.

\begin{definition}
    \textbf{Parity} is a property relating to (). If a system stays the same after a parity transformation, 
    it is said to have even parity. If it is inverted, it is said to have odd parity.
\end{definition}

\subsection{Differential Calculus}
 
\end{document}